\subsection{片段}

11.1.1. 设 E 是 Banach 空间。如果存在连续线性泛函 \(h \neq 0\) 和实数 k 使得 \(S = \{x|h(x) \leqslant k\}\)，则 E 的子集 S 称为闭半空间（参见 EVT, II, §2, no. 6）；此时 S 的边界是闭超平面 \(\{x|h(x) = k\}\)；它也称为 S 的边界，记作 \(\partial S\)。

11.1.2. 设 X 是类 \(C^{r}\) 流形，A 是 X 的一个子集。如果对每个 \(a \in A\)，都存在 X 在 a 处的图 \(c = (U, \varphi, E)\)，使得 \(\varphi(A \cap U)\) 是 E 的某个闭半空间的开子集，则称 A 是 X 的一个片段。置 \(\partial A = A - \mathring{A}\)（A 的非内部点的集合）。它是 A 的子流形，称为 A 的边界。\footnote{这一术语源于以下事实：A 自然地带有一个“带边界的流形”结构，其边界为 \(\partial A\)。关于这一范畴的定义，以及更一般的“带角流形”范畴，读者可参阅 H. Cartan, \emph{Séminaire 1961/62, Topologie Différentielle}, lectures 1–2–3（A. Douady），Benjamin, New York, 1969。需要指出的是，可以证明，每个边界为仿紧的“带边界流形”都同构于某个流形的闭片段。}若 \(a \in \partial A\)，则存在 X 在 a 处、以 a 为中心的图 \(c = (U, \varphi, E)\) 以及 E 的闭超平面 H，使得 \(\varphi(A \cap U)\) 是由 H 定义的 E 的某个闭半空间 \(S_{c}\) 中 0 的一个开邻域（EVT, II, § 2, no. 6），并且 \(\varphi(\partial A \cap U) = H \cap \varphi(A \cap U)\)。这样的图称为 A 在 a 处的适配图。于是，闭半空间 \(S_{c}\) 是使 \(\varphi(A \cap U)\) 成为开集的唯一 E 的闭半空间。

对于 X 的一个闭子集 A，它成为 X 的片段的充要条件是：\(\mathring{A}\) 在 A 中稠密，并且 \(A-\mathring{A}\) 在其每一点处都是 X 的余维 1 子流形。

11.1.3. 例子

\begin{enumerate}
\def\labelenumi{\alph{enumi})}
\item
  在 \(\mathbf{R}^{n}\) 中，半径 \(>0\) 的闭球是一个片段；其边界是相应的球面。
\item
  若 A 是流形 X 的一个片段，且 B 是 \(\partial A\) 的闭子集，则 A - B 是 X 的一个片段。
\item
  设 \(\varphi: X \to X'\) 是类 \(C^r\) 的流形态射，且 \(A'\) 是 \(X'\) 的一个片段。若 \(\varphi\) 与 \(\partial A'\) 横截（5.11.6），则 \(\varphi^{-1}(A')\) 是 \(X\) 的一个片段，其边界为 \(\varphi^{-1}(\partial A')\)。
\end{enumerate}

特别地，设 \(h: \mathbf{X} \to \mathbf{R}\) 是类 \(\mathbf{C}^r\) 的函数，且 \(a \in \mathbf{R}\)；假定不存在 \(x \in \mathbf{X}\) 使得 \(h(x) = a\) 且 \(d_x h = 0\)。集合 \(\{x | h(x) \leqslant a\}\) 是 \(\mathbf{X}\) 的一个闭片段，其边界为 \(h^{-1}(a)\)。

\begin{enumerate}
\def\labelenumi{\alph{enumi})}
\setcounter{enumi}{3}
\tightlist
\item
  若 \(r = \infty\) 且 X 局部紧，则 X 的每个紧子集都有一个由紧片段组成的邻域基本系。
\end{enumerate}

11.1.4. 设 X 是类 \(\mathbf{C}^{\mathrm{r}}\) 的流形，A 是 X 的一个片段。设 \(a \in \partial \mathrm{A}\)，\(v \in \mathrm{T}_a(\mathrm{X})\)。令 \(c = (\mathrm{U}, \varphi, \mathrm{E})\) 为 A 在 \(a\) 处的适配图（11.1.2），并令 \(h = \theta_c^{-1}(v)\) 为 E 中与 \(v\) 对应的元素（5.5.1）。令 \(\mathrm{S}_c\) 为 E 的闭半空间，使得 \(\varphi(\mathrm{A} \cap \mathrm{U})\) 是其开子集（11.1.2）。称 \(v\) 是 A 在 \(a\) 处的向内（相应地严格向内、向外、严格向外）向量，当且仅当 \(h\) 属于 \(\mathrm{S}_c\)（相应地属于 \(\dot{\mathrm{S}}_c\)、\(-\mathrm{S}_c\)、\(-\dot{\mathrm{S}}_c\)）；这个条件与适配图 \(c\) 的选择无关。我们把 \(\mathrm{T}_a^+(\mathrm{A})\)（相应地 \(\mathrm{T}_a^-(\mathrm{A})\)）记作 A 的向外（相应地向内）向量在 \(\mathrm{T}_a(\mathrm{X})\) 中的集合。这些是 \(\mathrm{T}_a(\mathrm{X})\) 的闭半空间，其边界包含 0。于是有

\[
\mathrm{T} _ {a} (\partial \mathrm{A}) = \mathrm{T} _ {a} ^ {+} (\mathrm{A}) \cap \mathrm{T} _ {a} ^ {-} (\mathrm{A}).
\]

